\chapter{第4章：黎曼面与模曲线 习题}

\difficultykey

\section{基础题 \easy}

\begin{problem}{商空间的局部坐标}
\easy
设 $\Gamma\subset \mathrm{SL}_2(\mathbf Z)$ 在上半平面 $\mathbb H$ 上作用离散。说明为什么在没有椭圆点的局部，商空间 $\Gamma\backslash\mathbb H$ 可以继承一维复流形结构。
\end{problem}
\begin{solution}
取一点 $\tau_0\in\mathbb H$，并假设其稳定子在 $\Gamma$ 中平凡。由于 $\Gamma$ 离散，可取一个足够小的圆盘 $U\ni\tau_0$，使得对任意 $\gamma\ne 1$，$\gamma U\cap U=\varnothing$。

于是投影映射
\[
\pi:U\longrightarrow \Gamma\backslash\mathbb H
\]
在像上是双全纯同构。换言之，$U$ 本身就给出了商空间的一张局部复坐标图。把这样的局部图拼起来，就得到 $\Gamma\backslash\mathbb H$ 在非椭圆点附近的一维复流形结构。
\end{solution}

\begin{problem}{识别二阶椭圆点}
\easy
证明 $i\in\mathbb H$ 是 $\mathrm{SL}_2(\mathbf Z)$ 作用下的二阶椭圆点。
\end{problem}
\begin{solution}
矩阵
\[
S=\begin{pmatrix}0&-1\\1&0\end{pmatrix}
\]
满足
\[
S(i)=\frac{-1}{i}=i.
\]
因此 $i$ 有非平凡稳定子。又
\[
S^2=-I,
\]
而在模群作用下 $-I$ 与恒等变换给出相同 Möbius 变换，所以由 $S$ 生成的稳定子在商意义下阶为 $2$。故 $i$ 是二阶椭圆点。
\end{solution}

\begin{problem}{识别三阶椭圆点}
\easy
令
\[
\rho=e^{2\pi i/3}=\frac{-1+\sqrt{-3}}{2}.
\]
证明 $\rho$ 是 $\mathrm{SL}_2(\mathbf Z)$ 作用下的三阶椭圆点。
\end{problem}
\begin{solution}
取
\[
ST=\begin{pmatrix}0&-1\\1&1\end{pmatrix}.
\]
直接计算得
\[
(ST)(\rho)=\frac{-1}{\rho+1}=\rho.
\]
所以 $\rho$ 有非平凡稳定子。又在 $\mathrm{PSL}_2(\mathbf Z)$ 中有
\[
(ST)^3=1,
\]
且 $ST$ 本身不是恒等变换，因此稳定子在商意义下阶为 $3$。故 $\rho$ 是三阶椭圆点。
\end{solution}

\begin{problem}{说明尖点的加入}
\easy
解释为什么给 $Y(1)=\mathrm{SL}_2(\mathbf Z)\backslash\mathbb H$ 加入无穷尖点以后，可以得到紧的模曲线 $X(1)$。
\end{problem}
\begin{solution}
在无穷远处引入局部参数
\[
q=e^{2\pi i\tau}.
\]
当 $\Im(\tau)\to\infty$ 时，$q\to0$，所以尖点附近可由穿心圆盘的商描述，并在 $q=0$ 处补上一个点。于是无穷尖点被填补成一个普通复点。

对 $\mathrm{SL}_2(\mathbf Z)$ 而言只有这一个尖点，因此加上它后，基本域的边界全部闭合，所得空间是紧的复曲线，这就是 $X(1)$。
\end{solution}

\begin{problem}{用模不变量参数化}
\easy
说明 $j$ 函数给出从 $X(1)$ 到黎曼球 $\mathbf P^1(\mathbf C)$ 的全纯参数化，因此 $X(1)$ 的亏格为零。
\end{problem}
\begin{solution}
$j(\tau)$ 是 $\mathrm{SL}_2(\mathbf Z)$ 不变的亚纯函数，所以它下降为 $X(1)$ 上的亚纯函数。它在尖点只有一个简单极点，并且在模曲线理论中给出全局坐标。因此
\[
j:X(1)\longrightarrow \mathbf P^1(\mathbf C)
\]
是双有理且实际上是双全纯同构。

既然 $X(1)$ 同构于黎曼球 $\mathbf P^1$，其亏格便是
\[
g(X(1))=0.
\]
\end{solution}

\begin{problem}{计算素数水平的指数}
\easy
证明对素数 $p$，有
\[
[\mathrm{SL}_2(\mathbf Z):\Gamma_0(p)]=p+1.
\]
\end{problem}
\begin{solution}
标准公式给出
\[
[\mathrm{SL}_2(\mathbf Z):\Gamma_0(N)]=N\prod_{\ell\mid N}\left(1+\frac1\ell\right).
\]
当 $N=p$ 为素数时，只有一个素因子 $p$，于是
\[
[\mathrm{SL}_2(\mathbf Z):\Gamma_0(p)]=p\left(1+\frac1p\right)=p+1.
\]
\end{solution}

\begin{problem}{计算十一级别亏格}
\easy
利用章节中的模曲线亏格公式
\[
g(X_0(N))=1+\frac{\mu_N}{12}-\frac{e_2(N)}{4}-\frac{e_3(N)}{3}-\frac{c_\infty(N)}{2}
\]
计算 $X_0(11)$ 的亏格。这里 $\mu_N=[\mathrm{SL}_2(\mathbf Z):\Gamma_0(N)]$，$e_2,e_3$ 分别是二阶与三阶椭圆点个数，$c_\infty$ 是尖点个数。
\end{problem}
\begin{solution}
对 $N=11$，有
\[
\mu_{11}=11\left(1+\frac1{11}\right)=12.
\]
因为 $11\equiv3\pmod4$，所以没有二阶椭圆点，即 $e_2(11)=0$。又 $11\equiv2\pmod3$，所以没有三阶椭圆点，即 $e_3(11)=0$。素数水平有两个尖点，故 $c_\infty(11)=2$。

代入公式得
\[
g(X_0(11))=1+\frac{12}{12}-0-0-\frac22=1.
\]
因此
\[
g(X_0(11))=1.
\]
\end{solution}

\begin{problem}{计算十四级别亏格}
\easy
用同一公式计算 $X_0(14)$ 的亏格。
\end{problem}
\begin{solution}
先算指数：
\[
\mu_{14}=14\left(1+\frac12\right)\left(1+\frac17\right)=14\cdot\frac32\cdot\frac87=24.
\]
因为 $14$ 的素因子中有 $7\equiv3\pmod4$，所以 $e_2(14)=0$。又因 $2\equiv2\pmod3$ 是其素因子，所以 $e_3(14)=0$。$14$ 平方因子自由，尖点个数为 $4$，故 $c_\infty(14)=4$。

于是
\[
g(X_0(14))=1+\frac{24}{12}-0-0-\frac42=1+2-2=1.
\]
所以
\[
g(X_0(14))=1.
\]
\end{solution}

\begin{problem}{由权二尖点形式得到微分}
\easy
设 $f\in S_2(\Gamma)$。证明
\[
\omega=f(\tau)\,d\tau
\]
在 $Y(\Gamma)$ 上下降为全纯微分，并在每个尖点处可延拓为全纯微分。
\end{problem}
\begin{solution}
先看群作用。若
\[
\gamma=\begin{pmatrix}a&b\\c&d\end{pmatrix}\in\Gamma,
\]
则权 $2$ 变换律给出
\[
f(\gamma\tau)=(c\tau+d)^2f(\tau).
\]
另一方面
\[
d(\gamma\tau)=\frac{d\tau}{(c\tau+d)^2}.
\]
相乘便得
\[
f(\gamma\tau)\,d(\gamma\tau)=f(\tau)\,d\tau.
\]
故它下降到商空间上。

再看尖点。取局部参数 $q=e^{2\pi i\tau/h}$，其中 $h$ 是该尖点宽度，则
\[
d\tau=\frac{h}{2\pi i}\frac{dq}{q}.
\]
因为 $f$ 是尖点形式，$f=\sum_{n\ge1}a_nq^n$，于是
\[
f(\tau)\,d\tau=\frac{h}{2\pi i}\left(\sum_{n\ge1}a_nq^{n-1}\right)dq,
\]
在 $q=0$ 处无极点，所以可全纯延拓到尖点。
\end{solution}

\begin{problem}{由微分反推权二形式}
\easy
设 $\omega$ 是 $X(\Gamma)$ 上的全纯微分。说明把它拉回到 $\mathbb H$ 后，可写成
\[
\omega=f(\tau)\,d\tau,
\]
其中 $f\in S_2(\Gamma)$。
\end{problem}
\begin{solution}
设 $\pi:\mathbb H\to Y(\Gamma)$ 是商映射。拉回后，$\pi^*\omega$ 是 $\mathbb H$ 上的全纯微分，因此必可写成
\[
\pi^*\omega=f(\tau)\,d\tau
\]
的形式，其中 $f$ 全纯。对任意 $\gamma\in\Gamma$，因为 $\omega$ 已下降到商空间，故
\[
\gamma^*(f(\tau)\,d\tau)=f(\tau)\,d\tau.
\]
这正等价于
\[
f(\gamma\tau)=(c\tau+d)^2f(\tau),
\]
即 $f$ 是权 $2$ 模形式。

再看尖点。若 $\omega$ 在紧化后的尖点处全纯，则用上题中的 $q$ 坐标反推可知 $f(\tau)$ 的 $q$-展开没有常数项，否则 $f(\tau)d\tau$ 会出现 $dq/q$ 型极点。故 $f$ 在每个尖点都消没，也就是
\[
f\in S_2(\Gamma).
\]
\end{solution}

\section{进阶题 \medium}

\begin{problem}{求素数水平的尖点宽度}
\medium
说明对 $\Gamma_0(p)$（$p$ 为素数），两个尖点可取为 $\infty$ 与 $0$，且它们的宽度分别为 $1$ 与 $p$。
\end{problem}
\begin{solution}
在 $\Gamma_0(p)$ 中，平移矩阵
\[
T=\begin{pmatrix}1&1\\0&1\end{pmatrix}
\]
属于该群，所以无穷尖点的最小正平移宽度为 $1$。

对尖点 $0$，用
\[
S=\begin{pmatrix}0&-1\\1&0\end{pmatrix}
\]
把它送到 $\infty$。于是其宽度等于最小正整数 $h$，使得
\[
S^{-1}\begin{pmatrix}1&h\\0&1\end{pmatrix}S=\begin{pmatrix}1&0\\-h&1\end{pmatrix}
\in\Gamma_0(p).
\]
这要求左下角被 $p$ 整除，即 $p\mid h$。最小正解是 $h=p$。故两个尖点的宽度分别为 $1$ 与 $p$。
\end{solution}

\begin{problem}{写出素数水平的亏格公式}
\medium
把章节中的一般亏格公式化为素数水平版本：证明
\[
g(X_0(p))=1+\frac{p+1}{12}-\frac{e_2(p)}{4}-\frac{e_3(p)}{3}-1.
\]
并将它简化为
\[
g(X_0(p))=\frac{p+1}{12}-\frac{e_2(p)}{4}-\frac{e_3(p)}{3}.
\]
\end{problem}
\begin{solution}
对素数 $p$，指数为
\[
\mu_p=p+1.
\]
尖点个数是 $2$，即 $c_\infty(p)=2$。把这些代入一般公式
\[
g(X_0(N))=1+\frac{\mu_N}{12}-\frac{e_2(N)}{4}-\frac{e_3(N)}{3}-\frac{c_\infty(N)}{2}
\]
便得到
\[
g(X_0(p))=1+\frac{p+1}{12}-\frac{e_2(p)}{4}-\frac{e_3(p)}{3}-1.
\]
首尾两个常数相消，于是化简为
\[
g(X_0(p))=\frac{p+1}{12}-\frac{e_2(p)}{4}-\frac{e_3(p)}{3}.
\]
这正是素数水平最常用的亏格公式。
\end{solution}

\begin{problem}{计算十七与二十三级别}
\medium
计算 $X_0(17)$ 与 $X_0(23)$ 的亏格。
\end{problem}
\begin{solution}
先算 $N=17$。因为 $17$ 是素数，
\[
g(X_0(17))=\frac{18}{12}-\frac{e_2(17)}{4}-\frac{e_3(17)}{3}.
\]
又 $17\equiv1\pmod4$，故 $e_2(17)=2$；$17\equiv2\pmod3$，故 $e_3(17)=0$。于是
\[
g(X_0(17))=\frac32-\frac24=1.
\]

再看 $N=23$。有
\[
g(X_0(23))=\frac{24}{12}-\frac{e_2(23)}{4}-\frac{e_3(23)}{3}.
\]
由于 $23\equiv3\pmod4$，$e_2(23)=0$；又 $23\equiv2\pmod3$，$e_3(23)=0$。故
\[
g(X_0(23))=2.
\]
所以
\[
g(X_0(17))=1,
\qquad
g(X_0(23))=2.
\]
\end{solution}

\begin{problem}{计算十五级别亏格}
\medium
计算 $X_0(15)$ 的亏格。
\end{problem}
\begin{solution}
指数为
\[
\mu_{15}=15\left(1+\frac13\right)\left(1+\frac15\right)=15\cdot\frac43\cdot\frac65=24.
\]
因为 $15$ 含有素因子 $3\equiv3\pmod4$，所以 $e_2(15)=0$。又它含有素因子 $5\equiv2\pmod3$，所以 $e_3(15)=0$。因为 $15$ 平方因子自由，所以尖点个数为 $4$。

代入公式：
\[
g(X_0(15))=1+\frac{24}{12}-0-0-\frac42=1+2-2=1.
\]
故
\[
g(X_0(15))=1.
\]
\end{solution}

\begin{problem}{由亏格得到维数}
\medium
利用同构
\[
H^0(X(\Gamma),\Omega^1)\cong S_2(\Gamma)
\]
以及 $\dim H^0(X,\Omega^1)=g(X)$，求出 $\dim S_2(\Gamma_0(11))$ 与 $\dim S_2(\Gamma_0(23))$。
\end{problem}
\begin{solution}
由前面计算，
\[
g(X_0(11))=1,
\qquad
g(X_0(23))=2.
\]
又对任意紧黎曼面 $X$，有
\[
\dim H^0(X,\Omega^1)=g(X).
\]
再结合
\[
H^0(X_0(N),\Omega^1)\cong S_2(\Gamma_0(N)),
\]
便得到
\[
\dim S_2(\Gamma_0(11))=1,
\qquad
\dim S_2(\Gamma_0(23))=2.
\]
\end{solution}

\begin{problem}{零亏格时没有微分}
\medium
证明若紧黎曼面 $X$ 的亏格为 $0$，则 $H^0(X,\Omega^1)=0$。
\end{problem}
\begin{solution}
一般定理给出
\[
\dim H^0(X,\Omega^1)=g(X).
\]
若 $g(X)=0$，则这个空间的维数就是 $0$，所以只能有零微分：
\[
H^0(X,\Omega^1)=0.
\]
应用到 $X(1)$ 上，就得到 $S_2(\mathrm{SL}_2(\mathbf Z))=0$。
\end{solution}

\begin{problem}{亏格一时没有零点}
\medium
设 $X$ 是亏格 $1$ 的紧黎曼面，$\omega\ne0$ 是全纯微分。证明 $\omega$ 没有零点。
\end{problem}
\begin{solution}
全纯微分的零点组成其除子 $(\omega)$，这是一个有效除子。其次数等于典范除子的次数：
\[
\deg(\omega)=\deg K_X=2g-2.
\]
若 $g=1$，则
\[
\deg K_X=0.
\]
一个有效除子的次数若为 $0$，只能是零除子。因此 $(\omega)=0$，也就是说 $\omega$ 没有零点。
\end{solution}

\begin{problem}{应用黎曼罗赫}
\medium
设 $X$ 是亏格 $g$ 的紧黎曼面，$D$ 是次数满足 $\deg D\ge 2g-1$ 的除子。用 Riemann--Roch 证明
\[
\ell(D)=\deg D+1-g.
\]
\end{problem}
\begin{solution}
Riemann--Roch 说
\[
\ell(D)-\ell(K-D)=\deg D+1-g.
\]
其中 $K$ 是典范除子，$\deg K=2g-2$。因此
\[
\deg(K-D)=(2g-2)-\deg D\le -1.
\]
次数为负的除子不可能有非零全纯截面，所以
\[
\ell(K-D)=0.
\]
代回 Riemann--Roch 即得
\[
\ell(D)=\deg D+1-g.
\]
\end{solution}

\begin{problem}{计算典范除子次数}
\medium
利用 Riemann--Roch 证明紧黎曼面 $X$ 的典范除子满足
\[
\deg K_X=2g-2.
\]
\end{problem}
\begin{solution}
对零除子 $D=0$ 使用 Riemann--Roch：
\[
\ell(0)-\ell(K)=1-g.
\]
这里 $\ell(0)=1$，因为只有常数函数。又 $\ell(K)=\dim H^0(X,\Omega^1)=g$。于是上式确实成立。

再把 $D=K$ 代入，得到
\[
\ell(K)-\ell(0)=\deg K+1-g.
\]
也就是
\[
g-1=\deg K+1-g.
\]
整理即得
\[
\deg K=2g-2.
\]
\end{solution}

\begin{problem}{构造受控极点函数}
\medium
设 $X$ 是亏格 $g$ 的紧黎曼面，$P\in X$。证明存在一个非常数亚纯函数，它在 $P$ 以外没有极点，并且在 $P$ 处的极点阶至多为 $g+1$。
\end{problem}
\begin{solution}
考虑除子
\[
D=(g+1)P.
\]
由 Riemann--Roch，
\[
\ell(D)-\ell(K-D)=\deg D+1-g=(g+1)+1-g=2.
\]
因此
\[
\ell(D)\ge2.
\]
这说明存在至少两个线性无关的函数属于 $L(D)$；其中一个可以取常数函数，所以还存在一个非常数函数 $f\in L(D)$。按定义，$f$ 的极点只可能出现在 $P$，且极点阶不超过 $g+1$。故命题成立。
\end{solution}

\begin{problem}{再次得到权二空间维数}
\medium
说明为什么对任意合同子群 $\Gamma$，都有
\[
\dim S_2(\Gamma)=g(X(\Gamma)).
\]
\end{problem}
\begin{solution}
章节中已经建立同构
\[
H^0(X(\Gamma),\Omega^1)\cong S_2(\Gamma).
\]
另一方面，对任意紧黎曼面 $X$，全纯微分空间的维数等于亏格：
\[
\dim H^0(X,\Omega^1)=g(X).
\]
把 $X$ 换成 $X(\Gamma)$，便得
\[
\dim S_2(\Gamma)=\dim H^0(X(\Gamma),\Omega^1)=g(X(\Gamma)).
\]
\end{solution}

\section{挑战题 \hard}

\begin{problem}{统一计算若干小水平}
\hard
分别计算 $X_0(11)$、$X_0(14)$、$X_0(15)$、$X_0(17)$、$X_0(23)$ 的亏格，并按亏格从小到大排列。
\end{problem}
\begin{solution}
前面已经分别算出
\[
g(X_0(11))=1,
\qquad
g(X_0(14))=1,
\qquad
g(X_0(15))=1,
\]
\[
g(X_0(17))=1,
\qquad
g(X_0(23))=2.
\]
因此从小到大排列为
\[
X_0(11),\ X_0(14),\ X_0(15),\ X_0(17)
\]
都具有亏格 $1$，而
\[
X_0(23)
\]
具有亏格 $2$，最大。
\end{solution}

\begin{problem}{计算三十七级别}
\hard
计算 $X_0(37)$ 的亏格，并由此求出 $\dim S_2(\Gamma_0(37))$。
\end{problem}
\begin{solution}
因为 $37$ 是素数，指数为
\[
\mu_{37}=38.
\]
又 $37\equiv1\pmod4$，故 $e_2(37)=2$；$37\equiv1\pmod3$，故 $e_3(37)=2$；尖点个数仍为 $2$。代入一般公式，得
\[
g(X_0(37))=1+\frac{38}{12}-\frac24-\frac23-\frac22.
\]
计算得
\[
g(X_0(37))=1+\frac{19}{6}-\frac12-\frac23-1=2.
\]
因此
\[
g(X_0(37))=2.
\]
再由
\[
\dim S_2(\Gamma_0(37))=g(X_0(37)),
\]
可知
\[
\dim S_2(\Gamma_0(37))=2.
\]
\end{solution}

\begin{problem}{零亏格模曲线上的结论}
\hard
设 $g(X_0(N))=0$。证明 $S_2(\Gamma_0(N))=0$，并解释这与 $X(1)$ 的情形一致。
\end{problem}
\begin{solution}
由章节同构，
\[
S_2(\Gamma_0(N))\cong H^0(X_0(N),\Omega^1).
\]
若 $g(X_0(N))=0$，则
\[
\dim H^0(X_0(N),\Omega^1)=0.
\]
因此
\[
S_2(\Gamma_0(N))=0.
\]
对 $X(1)$，我们已经知道它同构于 $\mathbf P^1$，故亏格为 $0$。于是
\[
S_2(\mathrm{SL}_2(\mathbf Z))=0,
\]
这与前面的基本事实完全一致。
\end{solution}

\begin{problem}{亏格一模曲线上的结论}
\hard
设 $g(X_0(N))=1$。证明 $S_2(\Gamma_0(N))$ 是一维的，并说明其中任一非零形式对应的全纯微分没有零点。
\end{problem}
\begin{solution}
因为
\[
\dim S_2(\Gamma_0(N))=g(X_0(N)),
\]
所以当 $g(X_0(N))=1$ 时，有
\[
\dim S_2(\Gamma_0(N))=1.
\]
任取非零 $f\in S_2(\Gamma_0(N))$，对应的全纯微分为
\[
\omega=f(\tau)\,d\tau.
\]
它是 $X_0(N)$ 上的非零全纯微分。由于亏格为 $1$，典范除子次数为
\[
2g-2=0.
\]
而 $(\omega)$ 是有效除子，其次数又等于 $0$，所以只能是零除子。故 $\omega$ 没有零点。
\end{solution}

\begin{problem}{描述局部分支模型}
\hard
说明在椭圆点附近，模曲线到 $X(1)$ 的局部模型分别可以写成 $z\mapsto z^2$ 或 $z\mapsto z^3$，并解释这对应二阶与三阶分支。
\end{problem}
\begin{solution}
若一点的稳定子在商意义下由有限循环群 $\mu_m$ 生成，则局部坐标可取成使该群作用为
\[
z\longmapsto \zeta_m z,
\qquad \zeta_m^m=1.
\]
在取商以后，不变量坐标是
\[
w=z^m.
\]
因此商映射的局部模型就是
\[
z\longmapsto z^m.
\]
当 $m=2$ 时得到二阶分支模型 $z\mapsto z^2$；当 $m=3$ 时得到三阶分支模型 $z\mapsto z^3$。这正对应于 $i$ 与 $\rho$ 处的二阶、三阶椭圆点。
\end{solution}

\begin{problem}{用黎曼罗赫得到球面}
\hard
利用 Riemann--Roch 证明：若紧黎曼面 $X$ 的亏格为 $0$，则 $X\cong \mathbf P^1(\mathbf C)$。
\end{problem}
\begin{solution}
任取一点 $P\in X$，考虑除子 $D=P$。由于 $g=0$，Riemann--Roch 给出
\[
\ell(P)-\ell(K-P)=\deg P+1=2.
\]
而
\[
\deg(K-P)=(-2)-1=-3<0,
\]
故 $\ell(K-P)=0$。于是
\[
\ell(P)=2.
\]
这说明存在非常数亚纯函数 $f\in L(P)$，它至多在 $P$ 处有一个简单极点，在其他地方无极点。

该函数定义了映射
\[
f:X\to \mathbf P^1.
\]
因为它只有一个简单极点，所以这个映射的次数为 $1$。次数为 $1$ 的非常数全纯映射必为双全纯同构。因此
\[
X\cong \mathbf P^1(\mathbf C).
\]
\end{solution}

\begin{problem}{从黎曼罗赫推出维数公式}
\hard
利用 Riemann--Roch 直接推出
\[
\dim H^0(X,\Omega^1)=g(X),
\]
再结合模曲线上的微分对应，得到
\[
\dim S_2(\Gamma)=g(X(\Gamma)).
\]
\end{problem}
\begin{solution}
设 $K$ 是紧黎曼面 $X$ 的典范除子。对零除子应用 Riemann--Roch：
\[
\ell(0)-\ell(K)=1-g.
\]
因为 $\ell(0)=1$，故
\[
1-\ell(K)=1-g,
\qquad \ell(K)=g.
\]
但 $\ell(K)$ 正是全纯微分空间的维数，即
\[
\dim H^0(X,\Omega^1)=g(X).
\]

现在取 $X=X(\Gamma)$。章节中已经证明
\[
H^0(X(\Gamma),\Omega^1)\cong S_2(\Gamma).
\]
于是
\[
\dim S_2(\Gamma)=\dim H^0(X(\Gamma),\Omega^1)=g(X(\Gamma)).
\]
这就得到权 $2$ 尖点形式空间的维数等于模曲线亏格。
\end{solution}

\begin{problem}{延拓到尖点的必要条件}
\hard
设 $f$ 是权 $2$ 模形式，且在某个尖点的局部参数 $q$ 下有展开
\[
f(\tau)=a_0+a_1q+a_2q^2+\cdots.
\]
证明微分 $f(\tau)d\tau$ 在该尖点全纯，当且仅当 $a_0=0$。
\end{problem}
\begin{solution}
设该尖点宽度为 $h$，局部参数为 $q=e^{2\pi i\tau/h}$。则
\[
d\tau=\frac{h}{2\pi i}\frac{dq}{q}.
\]
所以
\[
f(\tau)d\tau=\frac{h}{2\pi i}\left(a_0+a_1q+a_2q^2+\cdots\right)\frac{dq}{q}
\]
\[
=\frac{h}{2\pi i}\left(a_0q^{-1}+a_1+a_2q+\cdots\right)dq.
\]
可见在 $q=0$ 处是否有极点，完全由 $a_0q^{-1}dq$ 这一项决定。

因此，$f(\tau)d\tau$ 在尖点处全纯，当且仅当
\[
a_0=0.
\]
这正是权 $2$ 模形式成为尖点形式的条件。
\end{solution}
